% What the boundary algebra of p-adic AdS/CFT does not see
% Build: pdflatex boundary_blind (x3)
\documentclass[11pt]{article}
\usepackage{lmodern}
\usepackage[T1]{fontenc}
\usepackage{amsmath,amssymb,amsthm,mathtools}
\usepackage[margin=1.1in]{geometry}
\usepackage{booktabs}
\usepackage{microtype}
\usepackage[hidelinks]{hyperref}

\theoremstyle{plain}
\newtheorem{theorem}{Theorem}
\newtheorem{proposition}[theorem]{Proposition}
\newtheorem{corollary}[theorem]{Corollary}
\theoremstyle{remark}
\newtheorem{remark}[theorem]{Remark}

\newcommand{\Z}{\mathbb{Z}}
\newcommand{\C}{\mathbb{C}}
\newcommand{\F}{\mathbb{F}}
\newcommand{\PP}{\mathbb{P}}
\newcommand{\Jac}{\operatorname{Jac}}
\newcommand{\III}{\mathrm{III}}

\title{\bfseries What the boundary algebra of $p$-adic AdS/CFT does not see}
\author{Ruqing Chen\\[2pt]
\normalsize GUT Geoservice Inc., Rivi\`ere-Beaudette, Qu\'ebec J0P 1R0, Canada\\
\normalsize\texttt{ruqing@hotmail.com}}
\date{20 September 2026}

\begin{document}
\maketitle

\begin{abstract}
\noindent
In $p$-adic AdS/CFT the bulk is the Bruhat--Tits tree of $\mathrm{PGL}_2$ over a local
field and the boundary is $\PP^1(F)$. The boundary carries two standard operator algebras,
the crossed products $C(\PP^1(F))\rtimes\Gamma$ and $L^\infty(\PP^1(F))\rtimes\Gamma$, and
both have been completely determined: the first is a Cuntz--Krieger algebra with
$K_0=\Z^{r}\oplus\Z/(r-1)$ depending only on the rank of $\Gamma$, and the second is the
hyperfinite factor of type $\III_\lambda$ with $\lambda=q^{-2}$ or $q^{-1}$ according as
the quotient graph is bipartite or not (Robertson, 2005). Neither result has ever been
cited in the holography literature, and no $p$-adic AdS/CFT paper has used an
operator-algebraic invariant at all. We import them, and add one arithmetic layer. The
bulk Hecke--Laplacian channel has $K_0=\Z\oplus\Jac(Y)$, the critical group of the
quotient graph, whose order is a product of Frobenius traces
$\frac1n\prod_f(p+1-a_p(f))$. We exhibit two cubic graphs on eight vertices whose boundary
algebras agree in $K$-theory \emph{and} in von Neumann type, and whose bulk critical groups
have coprime orders $256=2^8$ and $363=3\cdot11^2$: the boundary data do not determine the
bulk arithmetic, not even which primes divide it. This is strictly stronger than
Robertson's own counterexample, which the type does separate. We do not conclude that
holography fails: Cornelissen and Marcolli proved that the boundary quantum statistical
mechanical system, with its Busemann time evolution, reconstructs the graph completely, so
what we exhibit is a precise gap in a hierarchy whose top is a reconstruction theorem.
Finally we observe that $\lambda\ne1$ is itself a diagnostic. Continuum holography is
forced to type $\III_1$ by theorems (mixing; half-sided modular inclusion); $p$-adic
holography is type $\III_\lambda$ with $\lambda<1$, so its modular flow is periodic and its
boundary algebra is not a Haag--Kastler algebra. All computations are exact and
machine-verified.
\end{abstract}

\section{Introduction}\label{sec:intro}

$p$-adic AdS/CFT \cite{GKPSW,HMSS} replaces anti-de Sitter space by the Bruhat--Tits tree
$X$ of $\mathrm{PGL}_2(F)$, $F$ a non-archimedean local field with residue field of size
$q$, and the conformal boundary by $\PP^1(F)$, the space of ends of $X$. A discrete
torsion-free cocompact lattice $\Gamma\le\mathrm{PGL}_2(F)$ --- free of finite rank $r$ ---
plays the role of a quotient identification, exactly as a Schottky group does for BTZ
\cite{HuangJepsen}, and $Y=\Gamma\backslash X$ is a finite $(q{+}1)$-regular graph on $n$
vertices. The field has matured on correlators, tensor networks, Wilson lines and finite
temperature \cite{GubserParikh,Marcolli20,HuangJepsen,Mondal}, but it has remained
kinematic in one specific respect: no invariant from operator algebras has ever been
computed in it.

That is unfortunate, because the relevant invariants were computed twenty years ago, in a
literature the holography community does not read. Robertson \cite{Rob05} determined both
boundary algebras of exactly this situation, and Cornelissen and Marcolli \cite{CM13}
determined precisely how much of the bulk each one remembers. The purpose of this note is
to import those results, to add one layer they do not contain, and to draw one physical
consequence.

\paragraph{What is imported.}
For $\Gamma$ a free uniform tree lattice of rank $r$, Robertson proved
\cite[Thm.~1]{Rob05} that $C(\partial X)\rtimes\Gamma$ is a Cuntz--Krieger algebra with
\begin{equation}\label{eq:K}
K_0\bigl(C(\partial X)\rtimes\Gamma\bigr)\cong\Z^{r}\oplus\Z/(r-1)\Z,
\qquad K_1\cong\Z^{r},
\end{equation}
$[1]$ generating the torsion summand --- so by Kirchberg--Phillips the algebra
\emph{depends only on $r$}. (The same $K$-groups were obtained independently in
\cite{CLM08}.) He proved further \cite[Thm.~2, Cor.~1]{Rob05} that
\begin{equation}\label{eq:type}
L^\infty(\partial X)\rtimes\Gamma\ \cong\ \text{hyperfinite }\III_\lambda,\qquad
\lambda=\begin{cases}q^{-2}&\Gamma\subset\mathrm{PSL}_2(F)\ (\,Y\ \text{bipartite}),\\
q^{-1}&\text{otherwise},\end{cases}
\end{equation}
the ratio set being generated by the Busemann cocycle; the rank-one case goes back to
Ramagge and Robertson \cite{RR97} and the general hyperbolic-group statement to
\cite{INO08,Bowen14}. Cornelissen and Marcolli \cite{CM13} then showed that the
$C^*$-level invariants ``reconstruct only topological information'' --- nine natural
equivalences all reduce to equality of the first Betti number --- while the boundary
\emph{quantum statistical mechanical} system, $C(\partial X)\rtimes\Gamma$ equipped with
the Busemann time evolution, reconstructs $Y$ up to isomorphism.

\paragraph{What is added.}
The bulk of this geometry has an invariant that none of these papers considers. The
Hecke--Laplacian $\Delta_p=(q{+}1)I-T_p$ of $Y$ has
\begin{equation}\label{eq:bulk}
\operatorname{coker}_\Z\Delta_p\cong\Z\oplus\Jac(Y),
\qquad
|\Jac(Y)|=\frac1n\prod_f\bigl(p+1-a_p(f)\bigr),
\end{equation}
where $\Jac(Y)$ is the critical (sandpile) group \cite{Biggs,BHN} and the product runs over
the weight-two newforms attached to $\Gamma$, so that $|\Jac(Y)|$ is a product of Frobenius
traces \cite[I]{ANCFT}. The critical group has appeared in physics --- as the global
symmetry group of fracton models on graphs \cite{GLS22,GLSS23} --- but never in holography.
Our contribution is Theorem~\ref{thm:main}: the boundary data of \eqref{eq:K} and
\eqref{eq:type} \emph{together} fail to determine \eqref{eq:bulk}, and fail badly enough
that two bulks in the same boundary class can have coprime arithmetic.

\paragraph{What we do not claim.}
Not that holography fails. \cite{CM13} is a reconstruction theorem: keep the time
evolution and the bulk comes back. What Theorem~\ref{thm:main} locates is a gap between two
rungs of a ladder whose top rung is a positive result. Nor do we claim \eqref{eq:K} or
\eqref{eq:type}; both are \cite{Rob05}.

\section{Setting}\label{sec:setting}

$Y=\Gamma\backslash X$ is a finite connected $(q{+}1)$-regular graph, $n=|V_Y|$,
$m=(q{+}1)n/2$, and
\[
r=\operatorname{rank}\Gamma=m-n+1,\qquad r-1=m-n=\tfrac{(q-1)n}{2}.
\]
The \emph{boundary channel} is the crossed product by the $\Gamma$-action on $\partial X$;
concretely $C(\partial X)\rtimes\Gamma\cong\mathcal O_{A^{\mathrm{nb}}}$ with
$A^{\mathrm{nb}}$ the non-backtracking matrix on the $2m$ oriented edges of $Y$
\cite{Rob05}, and $K_0=\operatorname{coker}(I-(A^{\mathrm{nb}})^{t})$,
$K_1=\ker(I-(A^{\mathrm{nb}})^{t})$ by \cite{Cuntz81}. The \emph{bulk channel} is
$\Delta_p$ on the $n$ vertices, with \eqref{eq:bulk}.

\begin{remark}[Two relations, two types]\label{rem:tworel}
Equation \eqref{eq:type} concerns the orbit relation of $\Gamma$ on $\partial X$. It must
be distinguished from the tail equivalence relation on the Cantor limit
$\varprojlim_kV_{Y_k}$ of a tower of covers, whose ratio set is all of $q^{\Z}\cup\{0\}$
and which therefore gives type $\III_{1/q}$ unconditionally \cite[III]{ANCFT}. The orbit
relation of a type-preserving lattice realises only $q^{2\Z}$, because the Busemann cocycle
of such an element has even value; this is the source of the dichotomy in \eqref{eq:type}
and it is easy to conflate the two.
\end{remark}

\begin{proposition}\label{prop:coarse}
The boundary $K$-theory \eqref{eq:K} is a function of $(q,n)$ alone, and the von Neumann
type \eqref{eq:type} is a function of $q$ and the bipartiteness of $Y$ alone. In
particular all $(q{+}1)$-regular graphs on $n$ vertices with the same bipartiteness have
the same boundary data.
\end{proposition}

\begin{proof}
$r-1=(q-1)n/2$ and $r=(q-1)n/2+1$ depend only on $(q,n)$; \eqref{eq:type} is as stated.
\end{proof}

\section{The bulk is not determined}\label{sec:main}

\begin{theorem}\label{thm:main}
There are two cubic graphs $Y_A,Y_B$ on eight vertices with
\begin{itemize}
\item the same boundary $K$-theory, $K_0=\Z^{5}\oplus\Z/4$ and $K_1=\Z^{5}$;
\item the same von Neumann type, both being non-bipartite and hence $\III_{1/q}$;
\item coprime bulk critical groups,
\[
\Jac(Y_A)=\Z/4\oplus\Z/4\oplus\Z/16,\quad |\Jac(Y_A)|=256=2^{8},
\]
\[
\Jac(Y_B)=\Z/33\oplus\Z/11,\quad\ \ |\Jac(Y_B)|=363=3\cdot11^{2}.
\]
\end{itemize}
Explicitly, with vertices $0,\dots,7$,
\[
\begin{aligned}
E(Y_A)&=\{03,06,07,12,15,16,24,25,34,37,45,67\},\\
E(Y_B)&=\{02,03,07,14,16,17,24,25,34,35,56,67\}.
\end{aligned}
\]
Consequently the boundary $K$-theory and the boundary von Neumann type, taken together, do
not determine $|\Jac(Y)|$, nor the set of primes dividing it, nor --- by \eqref{eq:bulk}
--- the Frobenius traces $a_p(f)$ of the associated automorphic forms.
\end{theorem}

\begin{proof}
Direct computation of the three invariants in exact integer arithmetic; this is check~(P)
of Section~\ref{sec:battery}. Both graphs contain an odd cycle, hence are non-bipartite.
\end{proof}

\begin{remark}[Comparison with \cite{Rob05}]\label{rem:compare}
Robertson gives his own counterexample to boundary rigidity \cite[Ex.~1.5]{Rob05}: the
theta graph and the dumbbell graph, both with $\pi_1=F_2$, non-isomorphic, with the same
$C(\partial X)\rtimes\Gamma$, realised concretely inside $\mathrm{PGL}_2(\mathbb Q_2)$. His
example is separated by the von Neumann type, $\III_{1/4}$ against $\III_{1/2}$ --- indeed
that is his point, that the measure-theoretic invariant is finer. Theorem~\ref{thm:main}
defeats both invariants at once, and does so in a way that is arithmetically meaningful
rather than merely graph-theoretic.
\end{remark}

Table~\ref{tab:named} shows the two channels on a list of standard quotients, and
Table~\ref{tab:spread} shows how much bulk arithmetic a single boundary class holds.

\begin{table}[ht]
\centering\small
\begin{tabular}{lrrrllr}
\toprule
$Y$ & $q$ & $n$ & $r$ & $K_0$(boundary) & vN type & $|\Jac(Y)|$\\
\midrule
$K_4$            & 2 &  4 &  3 & $\Z^3\oplus\Z/2$    & $\III_{1/q}$   & 16\\
$K_{3,3}$        & 2 &  6 &  4 & $\Z^4\oplus\Z/3$    & $\III_{1/q^2}$ & 81\\
prism $C_3\times K_2$ & 2 & 6 & 4 & $\Z^4\oplus\Z/3$ & $\III_{1/q}$   & 75\\
prism $C_4\times K_2$ & 2 & 8 & 5 & $\Z^5\oplus\Z/4$ & $\III_{1/q^2}$ & 384\\
M\"obius $M_8$   & 2 &  8 &  5 & $\Z^5\oplus\Z/4$    & $\III_{1/q}$   & 392\\
M\"obius $M_{10}$& 2 & 10 &  6 & $\Z^6\oplus\Z/5$    & $\III_{1/q^2}$ & 1815\\
prism $C_5\times K_2$ & 2 & 10 & 6 & $\Z^6\oplus\Z/5$ & $\III_{1/q}$  & 1805\\
Petersen         & 2 & 10 &  6 & $\Z^6\oplus\Z/5$    & $\III_{1/q}$   & 2000\\
Heawood          & 2 & 14 &  8 & $\Z^8\oplus\Z/7$    & $\III_{1/q^2}$ & 50421\\
$K_5$            & 3 &  5 &  6 & $\Z^6\oplus\Z/5$    & $\III_{1/q}$   & 125\\
$K_6$            & 4 &  6 & 10 & $\Z^{10}\oplus\Z/9$ & $\III_{1/q}$   & 1296\\
\bottomrule
\end{tabular}
\caption{The two channels. $K_1=\Z^r$ throughout. Within a block of equal $(q,n)$ the
boundary $K$-theory is constant; the type refines it by bipartiteness; $|\Jac(Y)|$ is
constrained by neither. Compare prism $C_5\times K_2$ and Petersen: same $(q,n)$, same
type, $|\Jac|=1805$ against $2000$. All entries exact.}
\label{tab:named}
\end{table}

\begin{table}[ht]
\centering\small
\begin{tabular}{rrlrr}
\toprule
$n$ & $r$ & $K_0$(boundary) & distinct $|\Jac|$ found & range\\
\midrule
 8 & 5 & $\Z^{5}\oplus\Z/4$ &  5 & $256$ -- $392$\\
10 & 6 & $\Z^{6}\oplus\Z/5$ & 16 & $960$ -- $2000$\\
12 & 7 & $\Z^{7}\oplus\Z/6$ & 38 & $2808$ -- $9747$\\
14 & 8 & $\Z^{8}\oplus\Z/7$ & 54 & $12100$ -- $46998$\\
\bottomrule
\end{tabular}
\caption{Bulk critical groups inside a single boundary $K$-theory class, from $60$ random
cubic quotients at each $n$. Sampled, not exhaustive: the true counts are at least these.}
\label{tab:spread}
\end{table}

\section{The type as a physical diagnostic}\label{sec:type}

The recent algebraic program in holography \cite{LL1,LL2,Witten22,CPW23,CLPW23} turns on
the fact that the large-$N$ algebra of a holographic CFT, and the algebra of a black hole
exterior, are of type $\III_1$, and that crossed products by the modular flow convert them
to semifinite type $\mathrm{II}$, which is what makes generalized entropy finite. Type
$\III_1$ is not an accident there: it is forced. If the correlation functions decay --- if
the algebra is generated by mixing operators in a KMS state --- the algebra must be
$\III_1$ \cite{FLMO23}; and the half-sided modular inclusion structure that
\cite{LL1,LL2} use forces $\III_1$ by Wiesbrock's theorem \cite{Wiesbrock93}. Local
algebras of relativistic quantum field theory are the hyperfinite $\III_1$ factor
\cite{BDF87}.

Equation \eqref{eq:type} says that $p$-adic holography is not in that class. Its boundary
algebra is $\III_\lambda$ with $\lambda=q^{-1}$ or $q^{-2}$, a Powers factor. The physical
content of $\lambda<1$ is that the modular flow is \emph{periodic}: the flow of weights is
the translation flow on $\mathbb R/T\mathbb Z$ with
\[
T=\frac{2\pi}{\log\lambda^{-1}}=\frac{2\pi}{\kappa\log q},\qquad\kappa\in\{1,2\},
\]
so that the modular Hamiltonian has spectrum in $(\log q^{\kappa})\Z$ rather than a
continuum. A $\III_\lambda$ algebra has a nontrivial discrete invariant --- Connes' $S$
invariant is $\lambda^{\Z}\cup\{0\}$ --- where a $\III_1$ algebra has none.

We draw one conclusion and refrain from a second. The conclusion: the type is a sharp,
computable obstruction to any claim that a $p$-adic holographic boundary theory is a local
quantum field theory in the Haag--Kastler sense, and it identifies precisely what would
have to change --- the residue field would have to be made continuous, $q\to1$ --- for the
$\III_1$ of \cite{LL1,Witten22} to be recovered. What we refrain from: we do not propose a
$p$-adic analogue of the crossed-product construction of \cite{Witten22}. For
$\III_\lambda$ with $\lambda<1$ the crossed product by the modular flow is a type
$\mathrm{II}_\infty$ algebra as well, but the relevant flow is periodic, and whether the
resulting trace has anything to do with an entropy is a question we have not investigated.

\section{Verification}\label{sec:battery}

The script \texttt{boundary\_blind.py} accompanying this note runs the checks of
Table~\ref{tab:battery} in $0.3$\,s, all passing; its output is archived as
\texttt{boundary\_blind\_output.txt}. Arithmetic is exact integer throughout except the two
rows marked \emph{sampled}.

\begin{table}[ht]
\centering\small
\begin{tabular}{clc}
\toprule
key & check & mode\\
\midrule
(K) & \eqref{eq:K} holds on 11 standard quotients, including $q=3,4$ & exact\\
(B) & $r-1=(q-1)n/2$: the boundary $K$-theory is a function of $(q,n)$ & exact\\
(J) & witness families at $(q,n)=(2,6),(2,8),(2,10)$: same boundary, different $\Jac$ & exact\\
(P) & Theorem~\ref{thm:main}: same $K$-theory, same type, coprime $|\Jac|$ & exact\\
(S) & Table~\ref{tab:spread}: the spread inside one boundary class & sampled\\
(R) & the Busemann cocycle is $q^{\Z}$-valued (normalisation check only) & exact\\
\bottomrule
\end{tabular}
\caption{The verification battery (\texttt{boundary\_blind.py}), $7/7$ passing.}
\label{tab:battery}
\end{table}

Row (R) deserves a caveat we state rather than bury: it verifies only that the cylinder
measures $\nu(\Omega_v)=q^{1-d(O,v)}$ give a $q^{\Z}$-valued Radon--Nikodym cocycle. That
the \emph{essential range} of that cocycle is attained --- which is what Krieger's theorem
\cite{Krieger76} needs --- is quoted from \cite{RR97,Rob05}, not verified here; no finite
computation could verify it.

\section{Scope}\label{sec:scope}

\emph{Quoted, not proved here.} Equations \eqref{eq:K} and \eqref{eq:type} and the
rigidity observation \cite{Rob05}; the reconstruction theorem \cite{CM13}; Krieger's
classification \cite{Krieger76}; the Cuntz machine \cite{Cuntz81}; the forcing of
$\III_1$ in the continuum \cite{FLMO23,Wiesbrock93,BDF87}.

\emph{New here.} Theorem~\ref{thm:main} and Tables~\ref{tab:named}--\ref{tab:spread}: the
arithmetic layer, and the observation that $K$-theory and type together are still blind to
it. The import of \cite{Rob05,CM13} into $p$-adic holography, where no operator-algebraic
invariant has previously been computed. The reading of $\lambda\ne1$ in
Section~\ref{sec:type} as a diagnostic separating $p$-adic from continuum holography.

\emph{Open.} (a) Whether the Cornelissen--Marcolli reconstruction \cite{CM13}, which
recovers $Y$ and hence $\Jac(Y)$, can be made to recover the arithmetic \emph{directly} ---
an explicit formula for $|\Jac(Y)|$ from the boundary KMS data would be a genuine
holographic dictionary entry for Frobenius traces. (b) The higher-rank case: for
$\widetilde A_n$ buildings the boundary type is $\III_{1/q}$ or $\III_{1/q^2}$ according to
the parity of $n$ \cite{CR03}, and holographic constructions on buildings exist
\cite{GMP22,Marcolli20}, but none is arithmetic; the bulk invariant there is the rank-two
critical group of \cite[X]{ANCFT}, and we have not computed the corresponding witnesses.
(c) Whether the periodicity of the modular flow has an entropic meaning.

\section*{Provenance}

Unrefereed preprint. The three exact computations behind Theorem~\ref{thm:main} --- the
Cuntz--Krieger $K$-groups, the bipartiteness, and the Smith normal form of the
Hecke--Laplacian --- are performed by the accompanying script and are reproducible in
seconds. Everything else in Sections~\ref{sec:intro}--\ref{sec:type} is quoted, with the
attribution set out in Section~\ref{sec:scope}. We searched the physics literature for any
prior use of an operator-algebraic invariant in $p$-adic AdS/CFT, and for any occurrence of
type $\III_\lambda$ with $\lambda\ne1$ in holography, and found none; a search cannot
establish a negative, and we state this as the limit of our knowledge.

\begin{thebibliography}{99}

\bibitem{BHN} R.~Bacher, P.~de~la~Harpe and T.~Nagnibeda,
\emph{The lattice of integral flows and the lattice of integral cuts on a finite graph},
Bull.\ Soc.\ Math.\ France \textbf{125} (1997), 167--198.

\bibitem{Biggs} N.~Biggs,
\emph{Chip-firing and the critical group of a graph},
J.\ Algebraic Combin.\ \textbf{9} (1999), 25--45.

\bibitem{Bowen14} L.~Bowen,
\emph{The type and stable type of the boundary of a Gromov hyperbolic group},
Geom.\ Dedicata \textbf{172} (2014), 363--386.

\bibitem{BDF87} D.~Buchholz, C.~D'Antoni and K.~Fredenhagen,
\emph{The universal structure of local algebras},
Comm.\ Math.\ Phys.\ \textbf{111} (1987), 123--135.

\bibitem{CPW23} V.~Chandrasekaran, G.~Penington and E.~Witten,
\emph{Large $N$ algebras and generalized entropy},
JHEP \textbf{04} (2023), 009.

\bibitem{CLPW23} V.~Chandrasekaran, R.~Longo, G.~Penington and E.~Witten,
\emph{An algebra of observables for de Sitter space},
JHEP \textbf{02} (2023), 082.

\bibitem{ANCFT} R.~Chen,
\emph{Algorithmic non-commutative class field theory, I--XVI}, preprints (2026);
here I (the two channels), III (the tail relation and its type) and X (rank two).

\bibitem{CLM08} G.~Cornelissen, O.~Lorscheid and M.~Marcolli,
\emph{On the $K$-theory of graph $C^*$-algebras},
Acta Appl.\ Math.\ \textbf{102} (2008), 57--69.

\bibitem{CM13} G.~Cornelissen and M.~Marcolli,
\emph{Graph reconstruction and quantum statistical mechanics},
J.\ Geom.\ Phys.\ \textbf{72} (2013), 110--117.

\bibitem{Cuntz81} J.~Cuntz,
\emph{A class of $C^*$-algebras and topological Markov chains II: reducible chains and the
Ext-functor for $C^*$-algebras}, Invent.\ Math.\ \textbf{63} (1981), 23--50.

\bibitem{CR03} P.~Cutting and G.~Robertson,
\emph{Type III actions on boundaries of $\widetilde A_n$ buildings},
J.\ Operator Theory \textbf{49} (2003), 25--44.

\bibitem{FLMO23} K.~Furuya, N.~Lashkari, M.~Moosa and S.~Ouseph,
\emph{Information loss, mixing and emergent type $\mathrm{III}_1$ factors},
JHEP \textbf{08} (2023), 111.

\bibitem{GMP22} E.~Gesteau, M.~Marcolli and S.~Parikh,
\emph{Holographic tensor networks from hyperbolic buildings},
JHEP \textbf{10} (2022), 169.

\bibitem{GLS22} P.~Gorantla, H.~T.~Lam and S.-H.~Shao,
\emph{Fractons on graphs and complexity},
Phys.\ Rev.\ B \textbf{106} (2022), 195139.

\bibitem{GLSS23} P.~Gorantla, H.~T.~Lam, N.~Seiberg and S.-H.~Shao,
\emph{Gapped lineon and fracton models on graphs},
Phys.\ Rev.\ B \textbf{107} (2023), 125121.

\bibitem{GKPSW} S.~S.~Gubser, J.~Knaute, S.~Parikh, A.~Samberg and P.~Witaszczyk,
\emph{$p$-adic AdS/CFT},
Comm.\ Math.\ Phys.\ \textbf{352} (2017), 1019--1059.

\bibitem{GubserParikh} S.~S.~Gubser and S.~Parikh,
\emph{Geodesic bulk diagrams on the Bruhat--Tits tree},
Phys.\ Rev.\ D \textbf{96} (2017), 066024.

\bibitem{HMSS} M.~Heydeman, M.~Marcolli, I.~Saberi and B.~Stoica,
\emph{Tensor networks, $p$-adic fields, and algebraic curves: arithmetic and the
$\mathrm{AdS}_3/\mathrm{CFT}_2$ correspondence},
Adv.\ Theor.\ Math.\ Phys.\ \textbf{22} (2018), 93--176.

\bibitem{HuangJepsen} A.~Huang and C.~B.~Jepsen,
\emph{Finite temperature at finite places},
JHEP \textbf{03} (2025), 097.

\bibitem{INO08} M.~Izumi, S.~Neshveyev and R.~Okayasu,
\emph{The ratio set of the harmonic measure of a random walk on a hyperbolic group},
Israel J.\ Math.\ \textbf{163} (2008), 285--316.

\bibitem{Krieger76} W.~Krieger,
\emph{On ergodic flows and the isomorphism of factors},
Math.\ Ann.\ \textbf{223} (1976), 19--70.

\bibitem{LL1} S.~Leutheusser and H.~Liu,
\emph{Causal connectability between quantum systems and the black hole interior in
holographic duality}, Phys.\ Rev.\ D \textbf{108} (2023), 086019.

\bibitem{LL2} S.~Leutheusser and H.~Liu,
\emph{Emergent times in holographic duality},
Phys.\ Rev.\ D \textbf{108} (2023), 086020.

\bibitem{Marcolli20} M.~Marcolli,
\emph{Holographic codes on Bruhat--Tits buildings and Drinfeld symmetric spaces},
Pure Appl.\ Math.\ Q.\ \textbf{16} (2020), 1--33.

\bibitem{Mondal} A.~Mondal, S.~Parikh, P.~Pradhan and R.~Sengar,
\emph{Toward holography on biregular trees},
Phys.\ Rev.\ D \textbf{112} (2025), 086011.

\bibitem{RR97} J.~Ramagge and G.~Robertson,
\emph{Factors from trees},
Proc.\ Amer.\ Math.\ Soc.\ \textbf{125} (1997), 2051--2055.

\bibitem{Rob05} G.~Robertson,
\emph{Boundary operator algebras for free uniform tree lattices},
Houston J.\ Math.\ \textbf{31} (2005), 913--935.

\bibitem{Wiesbrock93} H.-W.~Wiesbrock,
\emph{Half-sided modular inclusions of von Neumann algebras},
Comm.\ Math.\ Phys.\ \textbf{157} (1993), 83--92.

\bibitem{Witten22} E.~Witten,
\emph{Gravity and the crossed product},
JHEP \textbf{10} (2022), 008.

\end{thebibliography}

\end{document}
